\section{Modelos Híbridos}

Para problemas específicos e ao longo do desenvolvimento da teoria de filme fino nas últimas duas décadas, vários grupos desenvolveram modelos próprios, geralmente especializados em suas aplicações. Dentre os modelos desenvolvidos se aproximam do presente trabalho aqueles que usam metodologias hibridas ou acopladas, com o escoamento principal trocando informações com o filme. \citet{Lavalle2015} criou um acoplamento no qual o modelo reduzido se comunica com o escoamento principal através de condições de contorno, o gás enxerga o filme como uma superfície móvel com velocidade imposta, sendo que esta velocidade vem das equações do filme, enquanto do lado do filme o gás fornece a taxa de cisalhamento e a pressão local. Lavalle realiza o tratamento numérico do acoplamento, com esquema explícito entre as equações. Os resultados obtidos pelo modelo do mesmo são satisfatórios, da ordem de 5 \% quando comparado com DNS. Apesar de geometria utilizada neste trabalho ser móvel, o autor validou o modelo em geometrias simples, como canais bidimensionais.

Ainda dentro dos modelos híbridos, as formas de acoplamento dos modelos simplificados com o escoamento principal variam entre os autores. \citet{Qin2020} utiliza o modelo de lubrificação apresentado por \citet{Oron1997}, acoplado com um escoamento principal que resolve \textit{Level-set}, os dois escoamentos são acoplados a partir de uma viscosidade efetiva, que transmite os efeitos de lubrificação do filme para o gás e a partir termos de fonte de massa que garantem a conservação das propriedades entre o filme e o \textit{Level-set} na transição. 

\citet{Adechy2004} utiliza um acoplamento que transmite o efeito de arrasto de forma oriundo das ondas presentes e do formato do filme, além de termos fontes serem responsáveis pelo destacamento ou colisão de gotas entre os modelos. As alterações na geometria promovidas pelo escoamento são transferidos no formato de uma rugosidade efetiva para o escoamento principal.

Para escoamentos anulares em geometrias não triviais, como curvas, \citet{Zhang2024} utilizou um modelo de filme acoplado com métodos de rastreamentos de gotas, na comunicação direta do filme para o gás o autor utilizou a mesma rugosidade apresentada por \citet{Adechy2004}, contabilizando apenas os efeitos geométricos da ação do filme e não os efeitos na tensão cisalhante no escoamento principal.

A principal diferença que o modelo proposto possui em relação aos anteriores é a representação do domínio do filme de forma muito mais detalhada que as demais, a um custo intermediário entre as formulações simplificadas e o método multifásico completo. 